% ===============================================================
%  Chapter 9 -- Forced Writebacks on Dirty Sharing       (Part II)
% ===============================================================
\chapter{Forced Writebacks on Dirty Sharing}
\label{ch:dirty-sharing}

\begin{figure}[htbp]
\centering
\adjustbox{max width=\linewidth}{%
\begin{forest}
  chapter map,
  [Chapter 9
    [The transaction [Three messages]]
    [The root cause [S must be clean]]
    [Wasted write backs [Superseded writes]]
    [Sharing patterns [Four patterns] [Producer--consumer model]]
  ]
\end{forest}}
\caption[]{Structure of Chapter~9: one transaction, the
design rule behind it, and its cost in the sharing patterns where it occurs.}
\label{fig:ch9-map}
\end{figure}

% ---------------------------------------------------------------
\section{The Transaction}
\label{sec:dirty-transaction}
\sectionmap{[Transaction [Cache A holds M] [Cache B reads] [A sends data to B]
  [A also writes back] [Both end in S]]}
\index{MESI}\index{dirty sharing}%

In this report \emph{dirty sharing}\index{dirty sharing} means one situation: a
block is modified (dirty) in one cache and another cache reads it, so it becomes
shared. \Cref{fig:dirty-transaction} shows what MESI does. The holder sends the
data to the reader, writes the block back to memory, and both caches end in \stS.
The reader needs the data, so the transfer cannot be avoided. This chapter is
about the writeback.

\begin{figure}[htbp]
\centering
\begin{bigfig}[0.85]
\begin{tikzpicture}[font=\sffamily\scriptsize,
    ag/.style={draw, thick, rounded corners=2pt, minimum width=1.9cm,
               minimum height=0.6cm, align=center, font=\sffamily\small},
    msg/.style={-{Stealth}, thick}, st/.style={font=\sffamily\bfseries\small}]
  \node[ag, fill=tealC!15, draw=tealC] (A) at (0,0) {Cache A};
  \node[ag, fill=plumC!14, draw=plumC] (B) at (5.2,0) {Cache B};
  \node[ag, fill=amberC!16, draw=amberC] (M) at (2.6,-4.0) {Memory};
  \draw[dashed, coreGrey!60] (A) -- (0,-3.4);
  \draw[dashed, coreGrey!60] (B) -- (5.2,-3.4);
  \node[st, text=mesiBlue, anchor=east] at (-1.1,0) {\stM};
  \draw[msg, plumC] (5.2,-0.9) -- node[above]{1. readmiss (8 B)} (0,-0.9);
  \draw[msg, tealC] (0,-1.7) -- node[above]{2. send data (72 B)} (5.2,-1.7);
  \draw[msg, alertRed, very thick] (0,-2.4) -- node[right=6pt, pos=0.35, text=alertRed]{\textbf{3. write back (72 B)}} (2.6,-3.65);
  \node[st, text=mesiBlue] at (-0.8,-3.1) {\stS};
  \node[st, text=mesiBlue] at (6.0,-3.1) {\stS};
  \node[st, text=mesiBlue] at (6.5,0) {\stI};
  \node[text=coreGrey, align=left, anchor=west] at (7.0,-1.6)
       {necessary: B needs the data\\[6pt]\textcolor{alertRed}{questionable: does memory}\\\textcolor{alertRed}{need it now?}};
\end{tikzpicture}
\end{bigfig}
\caption[Dirty sharing in MESI]{Dirty sharing under MESI. Cache~A holds the block
in \stM\ and cache~B reads it. Message~2 is needed for correctness. Message~3 is
what this report calls a forced writeback.}
\label{fig:dirty-transaction}
\end{figure}

% ---------------------------------------------------------------
\section{The Root Cause: The Shared State Must Be Clean}
\label{sec:dirty-root}
\sectionmap{[Root cause [S blocks are evicted silently] [So S must match memory]
  [So M to S must update memory] [The write back follows]]}
\index{MESI}\index{MSI}\index{forced writeback}\index{coherence axioms}%

The prescribed textbook gives the reason when it describes how MSI answers a
\textit{readmiss} for a block in \stM:

\begin{quote}
``We assume that we can seamlessly evict a block in the S state. Hence, it is
necessary to write the data to the lower level before moving to the S
state.''~\cite[\S12.4.6, p.~602]{sarangi2025basic}
\end{quote}

MESI keeps this rule. \Cref{fig:root-chain} sets out the reasoning. A cache must
be able to drop an \stS\ block without telling anyone, which is what makes \stS\
cheap. A block can be dropped silently only if its data also survives somewhere
else, namely in memory. Every \stS\ block must therefore match memory, and a
block moving from \stM\ to \stS\ must update memory first.

\begin{figure}[htbp]
\centering
\begin{bigfig}[0.95]
\begin{tikzpicture}[font=\sffamily\scriptsize,
    s/.style={draw, thick, rounded corners=3pt, align=center, text width=2.55cm,
              minimum height=1.2cm},
    a/.style={-{Stealth[length=6pt]}, very thick, slateC}]
  \node[s, fill=tealC!12, draw=tealC] (r1) at (0,0) {S blocks may be\\evicted \textbf{silently}};
  \node[s, fill=plumC!12, draw=plumC] (r2) at (3.55,0) {so every S block\\must \textbf{match memory}};
  \node[s, fill=amberC!15, draw=amberC] (r3) at (7.1,0) {so a block moving\\M $\to$ S must make\\memory up to date};
  \node[s, fill=alertRed!12, draw=alertRed] (r4) at (10.65,0) {\textbf{forced write back}\\on every dirty\\sharing transaction};
  \draw[a] (r1) -- (r2); \draw[a] (r2) -- (r3); \draw[a] (r3) -- (r4);
\end{tikzpicture}
\end{bigfig}
\caption[Why MESI writes back on dirty sharing]{The reasoning behind the forced
writeback. The rule on the left makes shared blocks cheap to evict; the cost on
the right is paid on every dirty sharing transaction.}
\label{fig:root-chain}
\end{figure}

\begin{keybox}[The root cause]
Neither coherence axiom requires the \gls{fwb}\index{forced writeback}. It follows
from a design rule, that shared blocks are always clean, which was chosen so that
shared blocks can be evicted without a message.
\end{keybox}

No correct protocol can remove a cost that the axioms require, but a cost that
comes from a design rule can be removed by changing the rule, as long as the
axioms still hold. The rest of this report asks whether that change pays off.

% ---------------------------------------------------------------
\section{When the Writeback Is Wasted}
\label{sec:dirty-wasted}
\sectionmap{[Wasted write backs [Producer writes again] [Memory copy is out of date again]
  [Earlier write back was superseded] [Only the last one is needed]]}
\index{MESI}\index{forced writeback}%

A forced writeback is useful only if someone reads memory's copy before it goes
out of date. Take a producer core that keeps updating a block and a consumer core
that reads each new version (\cref{fig:wasted-timeline}). Under MESI every consumer
read forces a writeback, and every new producer write makes the copy just written
to memory out of date. Memory ends up with its final value only from the last
writeback before the block leaves the caches.

\begin{figure}[htbp]
\centering
\begin{bigfig}
\begin{tikzpicture}[font=\sffamily\scriptsize]
  \draw[very thick, coreGrey, -{Stealth}] (0,0) -- (13.4,0) node[right]{time};
  \foreach \x/\lab/\col in {0.8/{P writes\\v1}/tealC, 2.6/{C reads}/plumC,
                            4.4/{P writes\\v2}/tealC, 6.2/{C reads}/plumC,
                            8.0/{P writes\\v3}/tealC, 9.8/{C reads}/plumC,
                            11.9/{block evicted\\(silently)}/coreGrey} {
    \draw[\col, thick] (\x,-0.12) -- (\x,0.12);
    \node[above, align=center, text=\col] at (\x,0.15) {\lab};
  }
  \foreach \x/\v in {2.6/v1,6.2/v2,9.8/v3} {
    \draw[-{Stealth}, alertRed, thick] (\x,-0.15) -- (\x,-0.85);
    \node[draw, alertRed, rounded corners=1pt, fill=alertRed!8, below] at (\x,-0.85) {WB \v};
  }
  \node[text=leafC, align=center] at (11.9,-1.1) {memory already\\holds v3};
  \draw[decorate, decoration={brace, amplitude=4pt, mirror}] (2.1,-1.55) -- (10.3,-1.55)
    node[midway, below=5pt, text=alertRed, align=center]{three forced write backs; the copies of v1 and v2 in memory\\are overwritten before anything reads them};
\end{tikzpicture}
\end{bigfig}
\caption[Superseded writebacks]{A producer (P) and consumer (C) under MESI. Each
consumer read forces a writeback and the next producer write makes it out of date.
Only the writeback of v3 is still current when the block is evicted.}
\label{fig:wasted-timeline}
\end{figure}

Here two of the three forced writebacks are overwritten before memory's copy is
used. With $k$ rounds, $k-1$ of the $k$ forced writebacks are overwritten, and one
write to memory would have been enough.

% ---------------------------------------------------------------
\section{Sharing Patterns}
\label{sec:dirty-patterns}
\sectionmap{[Patterns [Private: never shared] [Read-only: never written]
  [Producer-consumer: one writer] [Migratory: writers take turns]]}
\index{MESI}\index{forced writeback}\index{sharing patterns!producer--consumer}%

Whether forced writebacks happen depends on how a program shares data.
Shared-memory programs are usually described by four sharing
patterns~\cite{culler1999parallel}\index{sharing patterns}. \Cref{fig:pattern-tree}
shows how to tell them apart and \cref{tab:patterns} gives the forced-writeback
cost of each under MESI.

\begin{figure}[htbp]
\centering
\adjustbox{max width=0.8\linewidth}{%
\begin{forest}
  decision tree,
  [Is the block used by\\more than one core?
    [Private, el=no]
    [Is it ever\\written?, el=yes
      [Read-only\\shared, el=no]
      [Do several cores\\write it in turn?, el=yes
        [Producer--consumer\\one writer and\\several readers, el=no]
        [Migratory\\each core reads\\then writes, el=yes]]]]
\end{forest}}
\caption[Decision tree: classifying a sharing pattern]{The four sharing patterns.
Only the two on the bottom row have a block that is modified in one cache and then
read by another.}
\label{fig:pattern-tree}
\end{figure}

\begin{table}[htbp]
\centering
\caption{The four sharing patterns and their forced-writeback cost under MESI.}
\label{tab:patterns}
\small
\begin{tabular}{@{}l p{4.6cm} p{4.4cm}@{}}
\toprule
Pattern & Example & Forced write backs (MESI) \\
\midrule
Private & a thread's own stack and locals & none: never shared \\
Read-only shared & an input matrix read by every thread & none: never modified \\
Producer--consumer & the \texttt{partial[]} results of \cref{lst:omp}, read
  by the combining thread & one per value produced \\
Migratory & a counter that each thread reads and then updates in turn &
  one per hand-off that begins with a read \\
\bottomrule
\end{tabular}
\end{table}

\label{sec:dirty-model}
So the size of the cost depends on the workload as much as on the protocol. A
program that mostly uses private and read-only data never pays it, and a pipeline
of producer and consumer threads pays it on every hand-off. For one producer and
$C$ consumers sharing a block under MESI in steady state, each round has the
producer's write invalidating the consumers (8\,B), the first consumer's read
missing, getting the data and forcing a writeback ($8+72+72=152$\,B), and each of
the other $C-1$ consumers missing and getting the data ($80$\,B each). The bytes
per round and the share of them that is the forced writeback are
\begin{equation}
  B(C) = 160 + 80\,(C-1), \qquad
  F(C) = \frac{72}{B(C)}
  \label{eq:pc-model}
\end{equation}

\begin{figure}[htbp]
\centering
\begin{tikzpicture}
  \begin{axis}[width=9.5cm, height=4.8cm, ybar stacked, bar width=20pt,
      xlabel={Number of consumers $C$}, ylabel={Bytes per value produced},
      symbolic x coords={1,2,4,8}, xtick=data, ymin=0, ymax=800,
      legend style={at={(0.03,0.97)}, anchor=north west, font=\scriptsize},
      grid=major, grid style={coreGrey!15},
      label style={font=\small}, tick label style={font=\scriptsize}]
    \addplot[fill=skyC!50, draw=skyC] coordinates {(1,88) (2,168) (4,328) (8,648)};
    \addlegendentry{control and data messages}
    \addplot[fill=alertRed!60, draw=alertRed,
             nodes near coords={\pgfmathprintnumber[fixed,precision=0]{\pgfplotspointmeta}\,\%},
             point meta=explicit, nodes near coords style={font=\scriptsize\bfseries, text=alertRed}]
      coordinates {(1,72) [45] (2,72) [30] (4,72) [18] (8,72) [10]};
    \addlegendentry{forced write back (labelled: share of total)}
  \end{axis}
\end{tikzpicture}
\caption[Forced writeback traffic in producer--consumer sharing]{\Cref{eq:pc-model}
for 1, 2, 4 and 8 consumers under MESI. The forced writeback is 72 bytes per round,
so its share is largest (45\,\%) when one consumer reads each value.}
\label{fig:pc-model}
\end{figure}

The share of traffic taken by the forced writeback therefore changes with the
workload (\cref{fig:pc-model}). It is largest in one-to-one producer--consumer
hand-offs and shrinks when many consumers read each value, so a study of the
problem has to vary the sharing pattern as well as the core count.

\label{sec:dirty-memory}
A forced writeback is paid for twice. It occupies the bus or network, which raises
the per-core demand $d$ of \cref{eq:util}, and it writes to the LLC or DRAM, using
memory bandwidth that all cores share and that does not grow with the core count.
When the writeback is later overwritten (\cref{sec:dirty-wasted}), neither cost
buys anything.
